% answers_short.tex —— 2006 年·普通高等学校招生考试·天津卷文科 · 简短答案速查
% 由 tools/make_short_answers.py 自 answers.tex 的 \gksolhead/\gkanswer 提取，
% 并与 2006/2006_数学真题_天津文_答案要点.md 逐题比对后生成，勿手改（改源 MD 后重跑）。

\begin{gknotice}
\textbf{使用说明}：本节只给最终答案便于快速核对，完整推导见解析版。答案为编写者独立推导，\textbf{非官方标准答案}。
\end{gknotice}

\gkpart{选择题}
\begin{center}\begin{tabular}{ccccccccccc}
\toprule
题号 & 1 & 2 & 3 & 4 & 5 & 6 & 7 & 8 & 9 & 10 \\
\midrule
答案 & A & B & B & A & C & D & B & D & D & B \\
\bottomrule
\end{tabular}\end{center}

\gkpart{填空题}
\begin{itemize}[leftmargin=2.6em]
  \item[\textbf{11.}] 35。
  \item[\textbf{12.}] $3/\sqrt{10}$。
  \item[\textbf{13.}] $3/4$。
  \item[\textbf{14.}] U（题面称“这两个圆”，但按同时与指定两条射线相切并要求切点位于射线上，当前条件只能确定一个候选；缺同版材料核查，不给出完整方程答案）。
  \item[\textbf{15.}] 20吨。
  \item[\textbf{16.}] 24。
\end{itemize}

\gkpart{解答题}
\begin{itemize}[leftmargin=2.6em]
  \item[\textbf{17.}] $\cos2\alpha=-3/5$，$\sin(2\alpha+\pi/4)=1/(5\sqrt2)$。
  \item[\textbf{18.}] (1) 0.243；(2) 0.995
  \item[\textbf{19.}] (1) 见解析版；(2) 见解析版
  \item[\textbf{20.}] (1) 无极值；(2) $\theta\in(\pi/3,\pi/2]$；(3) $a\in(-\infty,0]\cup[5/8,1)$
  \item[\textbf{21.}] (1) $\lambda=\pm1$；(2) 见解析版
  \item[\textbf{22.}] (1) $x^2-4y^2=1$；(2) 见解析版
\end{itemize}
